\documentclass[12pt]{article}
\usepackage[utf8]{inputenc}
\usepackage[T1]{fontenc}
\usepackage{hyperref}
\title{Bomitaba Language Resources for K-12 Curriculum Development}
\author{Compiled by Claude Sonnet 5.5 (claude-sonnet-5-5) for Andr\'as Kornai}
\date{}
\begin{document}
\maketitle

\section{Introduction}
Bomitaba (also Bamitaba, Leke, Mbomitaba; ISO 639-3 zmx; Glottocode bomi1238) is a Bantu language (Guthrie C14) spoken in northern Congo (Brazzaville), mainly along the Likouala-aux-Herbes river north of Epena, with a smaller community in the Central African Republic. Sources give 9,800 speakers (2000) and 7,224 (Global Recordings Network). There are Northern (Matoki) and Central (Epena) dialects. A 1988--89 survey reported that Bomitaba is used almost exclusively at home and in the village, that French is used in school and official contacts, and that Lingala is used in trade.

\section{Primary Educational Resources}
None found.

\section{Linguistic and Reference Materials}
\begin{itemize}
\item Wikipedia article: \url{https://en.wikipedia.org/wiki/Bomitaba_language}
\item Gardner's 1990 University of North Dakota master's thesis, a sociolinguistic survey of the Bomitaba: \url{https://commons.und.edu/theses/5067}
\item Tyndale article, ``Language use in the Epena District of northern Congo'': \url{https://www.tyndale.ca/publications/language-use-in-the-epena-district-northern-congo}
\item Another search result, whose content was not examined: \url{https://core.ac.uk/works/153732631}
\end{itemize}

\section{Religious and Community Resources}
Joshua Project has a people-group profile, which gives demographic and religious background but contains no language materials: \url{https://www.joshuaproject.net/index.php/people_groups/10927/CF}

\section{Digital and Audio Resources}
Global Recordings Network lists Bomitaba recordings: \url{https://globalrecordings.net/en/language/zmx}. The contents were not examined.

\section{Curriculum Development Context}
The survey evidence suggests the language is still used in homes and villages, while schooling is in French. No Bomitaba school programme, primer or dictionary was found. A curriculum would need new materials and an orthography developed with the community.
\end{document}